\bookchapter{Глава 48: Знакомство с Android — от операционной системы до экосистемы разработки}

\chapterintro{Поздравляем с этим решением — из «человека, который пользуется телефоном» стать «человеком, который делает приложения». Прежде чем браться за дело, расстелим карту: откуда взялся Android? Почему разработчики всего мира пишут для него приложения? Что происходит между моментом, когда вы стучите по клавиатуре, и появлением приложения на телефоне пользователя?}

\booksection{Происхождение Android: от «маленького робота» до «крупнейшей системы мира»}

История Android началась в 2003 году: инженер Энди Рубин основал компанию Android Inc. с простой целью — создать умную операционную систему, «лучше понимающую камеры». В 2005 году её купил Google, в 2007-м вместе с несколькими десятками производителей был создан альянс Open Handset Alliance (OHA), а в сентябре 2008 года официально поступил в продажу первый Android-смартфон — HTC Dream (G1).

Дальнейшее вы знаете: к 2026 году Android — мобильная операционная система с самой большой долей рынка в мире, а активных устройств — миллиарды. Ваш телефон, ТВ-приставка, экран в машине, умные часы — почти наверняка работают на Android. Что это значит? Что каждая написанная вами строка кода может дойти до самого широкого круга людей на планете.

\begin{tcbinfo}{Совет: открытые корни Android}
Ядро Android (AOSP, Android Open Source Project) открыто: любой производитель может взять его и переделать под себя. Поэтому у Xiaomi есть MIUI (теперь HyperOS), у Huawei — слой совместимости HarmonyOS, у Samsung — One UI. Для разработчика это благо: вы изучаете «разработку под Android», а не «разработку под конкретный бренд» — одно умение работает на всех.

\end{tcbinfo}

Забавна и эволюция номеров версий: начиная с 1.5 Cupcake («капкейк») каждая крупная версия носила имя десерта — пончик, эклер, замороженный йогурт, пряничный человечек… И только Android 10 в 2019 году перешла на сухие цифры. На август 2026 года последняя официальная версия — Android 17, а «минимально поддерживаемая версия», которая действительно важна при разработке, у массовых приложений обычно между Android 8 (API 26) и Android 10 (API 29): поддерживая устройства последних нескольких лет, вы охватываете достаточно большую аудиторию.

\begin{center}
\begin{tabularx}{\linewidth}{|X|X|X|}
\hline
\textbf{Версия} & \textbf{Кодовое имя / год} & \textbf{Важное для разработчика} \tabularnewline
\hline
Android 8 (API 26) & Oreo · 2017 & Каналы уведомлений, ограничения фоновой работы — нижняя планка для многих старых устройств \tabularnewline
Android 10 (API 29) & 2019 & Появилось изолированное хранилище (scoped storage), правила доступа к файлам сильно изменились \tabularnewline
Android 12 (API 31) & 2021 & Динамические цвета Material You — большой шаг вперёд в облике интерфейса \tabularnewline
Android 13 (API 33) & 2022 & Разрешение на уведомления в рантайме: система спрашивает, «разрешать ли присылать уведомления» \tabularnewline
Android 14-15 (API 34-35) & 2023-2024 & Прогнозируемое возвращение, страницы по 16 КБ, дальнейшее ужесточение приватности \tabularnewline
Android 16-17 (API 36-37) & 2025-2026 & Более частые минорные релизы, интеграция ИИ-возможностей (например, расширения Gemini) \tabularnewline
\hline
\end{tabularx}
\end{center}

\begin{center}\small\color{gray}Краткая история версий Android (на август 2026 года; главные вехи для разработки)\end{center}

\booksection{На каком языке писать: Java или Kotlin?}

За историю разработки под Android сменились два основных языка. Java — «ветеран-основатель»: она была с Android с самого рождения. Kotlin — объявленный Google в 2017 году «гражданином первого сорта», а с 2019 года официальная линия — «Kotlin-first». К 2026 году новые проекты почти все пишут на Kotlin, а старые постепенно мигрируют.

Скажем честно: если вы совсем новичок, учите сразу Kotlin и не оглядывайтесь на Java. Причины три. Первая: кода на Kotlin на 30–40\% меньше, чем на Java (та же функциональность — короче и яснее). Вторая: Kotlin от природы защищает от «нулевых указателей» — кошмара новичков. Третья: все современные уроки, официальная документация и ИИ-инструменты для программирования (помните нашу книгу о программировании с ИИ?) ориентированы на Kotlin. Java встретится лишь при сопровождении старых проектов и чтении старого кода.

Стиль Java (многословный)

\begin{lstlisting}
public class MainActivity extends Activity \{
@Override
protected void onCreate(Bundle b) \{
super.onCreate(b);
setContentView(R.layout.activity);
\}
\}
\end{lstlisting}

Стиль Kotlin (лаконичный)

\begin{lstlisting}
class MainActivity : ComponentActivity() \{
override fun onCreate(b: Bundle?) \{
super.onCreate(b)
setContent \{ App() \}
\}
\}
\end{lstlisting}

Одно и то же действие, а в Kotlin минус куча фигурных скобок и шаблонных строк. В главе 50 мы займёмся им всерьёз, а пока запомните: «общий язык» разработки под Android называется Kotlin.

\booksection{Как рождается приложение?}

Многие думают, что «разработка приложения» — это сесть за компьютер и лихорадочно писать код. На самом деле от идеи до телефона пользователя приложение проходит примерно такой конвейер:

Идея

Проектирование (интерфейс + функции)

Пишем код

Сборка (APK/AAB)

Тесты (эмулятор + реальное устройство)

Публикация (магазин/установщик)

\begin{center}\small\color{gray}Рисунок 1: полный путь Android-приложения — от идеи до полки магазина\end{center}

В этом процессе есть два понятия, которые нужно усвоить прямо сейчас:

\begin{itemize}
  \item APK: формат установочного пакета Android-приложения (.apk) — как .exe в Windows: его можно поставить на телефон. Годится для собственных тестов и «перебросить другу поиграть».
  \item AAB: Android App Bundle (.aab) — «продвинутый формат упаковки» для публикации в Google Play. Магазин динамически подрезает пакет под модель телефона пользователя, чтобы каждый скачивал меньше. Самому установить его нельзя — только через магазин.
\end{itemize}

И ещё один термин, который будет звучать постоянно: SDK (Software Development Kit — комплект средств разработки). Это весь набор «экипировки» от Android: API-библиотеки разных версий, отладочные инструменты, образы эмулятора и т. д. Классы вроде Activity и Button, которые вы вызываете в коде, приходят именно из SDK.

\booksection{Панорама технологий Android: ваш «список вооружений»}

Перед началом обучения познакомьтесь со стандартным набором типичного Android-проекта 2026 года. Если сейчас список непонятен — не страшно: перечитаете его после всей книги, и всё встанет на свои места:

\begin{center}
\begin{tabularx}{\linewidth}{|X|X|X|}
\hline
\textbf{Слой} & \textbf{Основной выбор (2026)} & \textbf{Для чего нужен} \tabularnewline
\hline
Язык & Kotlin 2.x & Язык написания логики \tabularnewline
Интерфейс & Jetpack Compose + Material 3 & Интерфейс рисуется кодом (главная линия книги) \tabularnewline
Архитектура & MVVM + ViewModel & Разделяет интерфейс и логику — код не превращается в кашу \tabularnewline
Хранение данных & Room (база данных) + DataStore (настройки) & Сохраняет данные на телефоне \tabularnewline
Сеть & Retrofit + OkHttp + корутины Kotlin & Общение с сервером \tabularnewline
Инструмент сборки & Gradle + Kotlin DSL & Превращает код в установочный пакет \tabularnewline
Среда разработки & Android Studio & Официальная IDE-«всё в одном» от Google \tabularnewline
\hline
\end{tabularx}
\end{center}

\begin{center}\small\color{gray}Основной стек Android-разработки 2026 года одним взглядом\end{center}

Не пугайтесь этого набора английских слов. Он как кастрюли и половники на кухне: сейчас достаточно знать, что «готовить надо на кухне», а как пользоваться каждым инструментом — разберём по главам, один за другим. И хорошая новость: при создании нового проекта Android Studio сама настроит весь этот набор — вам останется только «жарить».

\booksection{Android против iOS: первый «урок выбора стороны» для разработчика}

Возможно, у вас в руках и Android, и устройство Apple. Как разработчику стоит заранее знать разницу в опыте с двух сторон — это не «какая платформа лучше», а трезвый расчёт, «с какой стороны вам удобнее начать»:

\begin{center}
\begin{tabularx}{\linewidth}{|X|X|X|}
\hline
\textbf{Критерий} & \textbf{Разработка под Android} & \textbf{Разработка под iOS} \tabularnewline
\hline
Компьютер для разработки & Подходит Windows / Mac / Linux & Только Mac (Xcode эксклюзивен) \tabularnewline
Инструмент разработки & Android Studio (бесплатно) & Xcode (бесплатно) + 99 долларов в год за публикацию \tabularnewline
Проверка при публикации & Мягкая: от нескольких часов до пары дней & Строгая: 1–7 дней — обычное дело \tabularnewline
Установка на устройство и тесты & APK распространяется и ставится свободно & Распространение ограничено: для тестов нужны подпись / TestFlight \tabularnewline
Условия заработка & В основном реклама / локальные магазины & В App Store привычка платить за внутренние покупки лучше \tabularnewline
Стоимость обучения & Понятий больше (четыре главных компонента и т. д.) & Экосистема закрытая, но целостная, документация отличная \tabularnewline
\hline
\end{tabularx}
\end{center}

\begin{center}\small\color{gray}Сравнение путей разработки под Android и iOS (август 2026 года)\end{center}

Скажем честно: если ваш компьютер не Mac, выбирать не из чего — Android единственный вариант. Если Mac есть, восьмую и девятую части этой книги можно читать параллельно: логика программирования общая, и после освоения одной стороны вторая пойдёт куда быстрее.

\begin{tcbpractice}{Попробуйте сами: посмотрите на «Android» в своём телефоне}
Откройте в телефоне [Настройки] [О телефоне] (или «Об этом устройстве») и найдите строку «Версия Android» — нажмите на неё несколько раз подряд, и сработает пасхалка: мини-игра от системы. Запишите номер своей версии и сверьтесь с таблицей выше: к какому поколению относится ваш телефон. Ваш телефон и станет первым «тестовым устройством» для разработки.

\end{tcbpractice}

\begin{tcbsummary}{Краткие итоги главы}
\begin{itemize}
  \item Android родился в 2008 году; к 2026-му — мобильная система с самой большой долей рынка в мире, миллиарды активных устройств
  \item Последняя официальная версия — Android 17; массовые приложения поддерживают минимум Android 8–10 — так больше пользователей
  \item Новые проекты — всегда на Kotlin: кода на 30\%+ меньше, защита от нулевых указателей, официальный выбор
  \item APK — для установки самому, AAB — для магазина; SDK — официальный набор для разработки
  \item Стандартный стек 2026: Kotlin + Compose + MVVM + Room + Retrofit + Gradle + Android Studio
\end{itemize}

\end{tcbsummary}
